The released pushpin branch, paired RGB manipulations and 138 held-out original photos were established in the preceding audit. The inspection endpoints and paired target-policy tests are follow-ups on those original photos. The V3 protocol was frozen before examining new scores. A separate design note added upstream-only fine-tuning after development revealed that the class-wide rule-aware head failed the cell ranking competence check. That note precedes inspection of held-out V3 queue results. The juice eligibility and edit lock is separate and precedes scoring there. We keep all registered policies, seeds and interfaces.

\paragraph{Scene and spatial cutoffs.}
Each scene cutoff is the 63rd smallest maximum of 69 normal validation maps, with an alert requiring a strictly greater score. The spatial cutoff at recall $r$ is the $\lceil 138r\rceil$th largest of the development target scores on scenes that alert. Both vacated and recipient cells are included, even when one is difficult to localize. The spatial decision is inclusive. With fewer than $\lceil 138r\rceil$ targets available after gating, the queue cutoff is set to zero and feasibility is reported. The fixed-cost cutoff is the $(\lfloor0.5 n_{\rm dev}\rfloor+1)$th largest score on untouched cells of gated development faults, with a strict spatial comparison. This gives at most 0.5 untouched cells per development fault. If fewer negatives remain after gating, the cutoff is zero. On valid shams and sources, any queued cell is a false mark. For the vacancy-only variant, the vacated cell is the sole positive. Rank results use the raw maps before scene gating. Exact ties break by cell index. The $2/3$ boundary tie rate is archived.

\paragraph{Decision criteria.}
The primary endpoint is the difference in untouched marks per extracted fault, footprint minus class-wide. A useful difference is at least one fewer mark with a 95\% interval excluding zero, 90\% or higher achieved target recall with at most a five-point decline, at least 95\% fault and vacancy scene recall with at most a three-point decline, and no more than five extra points of valid-sham alerts. A mechanistic interpretation requires class-wide development scene recall of at least 90\% for both fault types, fault-over-sham wins of at least 90\%, and a cell AUROC of at least 90\%. An explanation of released behavior additionally requires broad untouched-cell change under the class-wide control, above vacancy and sham. The seed and photo-half checks keep the sign of the primary contrast. All intervals condition on the fitted heads and cutoffs.

\paragraph{Regional analogue and diagnostic response.}
The controlled low-FPR region measure uses binary saturation, so each compartment counts once when its score crosses a shared cutoff. We include the two target regions and thirteen untouched regions per fault plus 30 negative regions from that original and its sham. For the fault-only curve we use 13 negatives per fault. We use linear interpolation through ties and integrate up to a 0.05 pooled negative-cell rate, then divide the area by 0.05. Both curves are descriptive sweeps across the test set. The strict scene and spatial cutoffs instead come from development. Bootstrap replicates resample complete photo groups and recalculate the pooled sweep. The official pixel AU-sPRO allows area-based saturation inside ambiguous native regions \cite{bergmann2022}. Our compartments have an exact action label and are indivisible review marks.

A cell containing one big change and one with broad changes both count in the any-pixel summary. The complementary area fraction counts pixels with an absolute shift of at least 0.5 on the same 13 cells, weighting by cell area within each original. We report per-cell upper quantiles and means across cells. None of those source comparisons enters stand-alone inspection.

